\documentclass[11pt,a4paper]{article}
\usepackage[utf8]{inputenc}
\usepackage[T1]{fontenc}
\usepackage[spanish,es-noshorthands]{babel}
\usepackage[margin=2.5cm]{geometry}
\usepackage{amsmath,amssymb,mathtools}
\usepackage{enumitem}
\usepackage{booktabs}
\usepackage{tikz}
\usepackage{pgfplots}
\pgfplotsset{compat=1.18}

\newcounter{ejnum}
\newenvironment{ejercicio}{\refstepcounter{ejnum}\par\medskip\noindent\textbf{Ejercicio \theejnum.}~}{\par}
\newenvironment{solucion}{\par\noindent\textit{Solución.}~}{\par\vspace{1em}}

\title{Práctica 2: Covarianza y función generatriz de momentos \\ \large Unidad 4: Función generadora de momentos}
\author{Probabilidad y Estadística (5765) \\ Universidad Nacional del Sur}
\date{}

\begin{document}
\maketitle

\begin{ejercicio}
Sean $X, Y$ variables aleatorias discretas con función de probabilidad conjunta dada por la tabla:
\[
\begin{array}{c|ccc}
 & Y=0 & Y=1 & Y=2 \\ \hline
X=0 & 0.20 & 0.10 & 0.00 \\
X=1 & 0.10 & 0.30 & 0.30
\end{array}
\]
Calcular $\text{Cov}(X,Y)$ y el coeficiente de correlación $\rho_{X,Y}$.
\end{ejercicio}

\begin{solucion}
Marginales: $P(X=0)=0.30$, $P(X=1)=0.70$; $P(Y=0)=0.30$, $P(Y=1)=0.40$, $P(Y=2)=0.30$.

\medskip
\textbf{Esperanzas:}
\[
E(X) = 0(0.30)+1(0.70) = 0.70, \qquad E(Y) = 0(0.30)+1(0.40)+2(0.30) = 1.00.
\]

\textbf{Momentos de segundo orden:}
\[
E(X^2) = 0^2(0.30)+1^2(0.70) = 0.70 \ \Rightarrow\ \text{Var}(X) = 0.70-0.70^2 = 0.70-0.49=0.21.
\]
\[
E(Y^2) = 0^2(0.30)+1^2(0.40)+2^2(0.30) = 0.40+1.20 = 1.60 \ \Rightarrow\ \text{Var}(Y) = 1.60-1.00^2 = 0.60.
\]

\textbf{Esperanza del producto:}
\[
E(XY) = \sum_{x,y} xy\, p(x,y) = 0+0+0+0 + (1)(1)(0.30)+(1)(2)(0.30) = 0.30+0.60=0.90.
\]

\textbf{Covarianza:}
\[
\text{Cov}(X,Y) = E(XY)-E(X)E(Y) = 0.90 - (0.70)(1.00) = 0.20.
\]

\textbf{Correlación:}
\[
\sigma_X = \sqrt{0.21}\approx 0.4583, \qquad \sigma_Y=\sqrt{0.60}\approx 0.7746.
\]
\[
\rho_{X,Y} = \frac{0.20}{(0.4583)(0.7746)} \approx \frac{0.20}{0.3550} \approx 0.5634.
\]
Existe una asociación lineal positiva moderada entre $X$ e $Y$.
\end{solucion}

\begin{ejercicio}
Sea $X$ el número de horas de estudio semanales de un estudiante e $Y=3X+7$ el puntaje que obtiene en un examen estandarizado (relación lineal exacta según el modelo simplificado del ejercicio). Si $\text{Var}(X)=4$, calcular $\text{Cov}(X,Y)$ y $\rho_{X,Y}$.
\end{ejercicio}

\begin{solucion}
Usando bilinealidad de la covarianza: $\text{Cov}(X,Y) = \text{Cov}(X, 3X+7) = 3\,\text{Cov}(X,X) + \text{Cov}(X,7)$.

Como $\text{Cov}(X,X)=\text{Var}(X)=4$ y la covarianza con una constante es $0$:
\[
\text{Cov}(X,Y) = 3(4) = 12.
\]
Además, $\text{Var}(Y) = \text{Var}(3X+7) = 3^2\text{Var}(X) = 9(4)=36$, por lo tanto $\sigma_Y=6$, $\sigma_X=2$.
\[
\rho_{X,Y} = \frac{12}{(2)(6)} = \frac{12}{12} = 1.
\]
Tiene sentido: al ser $Y$ una función lineal creciente exacta de $X$, la correlación es perfecta ($\rho=1$).
\end{solucion}

\begin{ejercicio}
Se invierte en dos activos financieros $X$ (acciones) e $Y$ (bonos), con $\text{Var}(X)=100$, $\text{Var}(Y)=25$ y $\text{Cov}(X,Y)=-30$ (correlación negativa, típica de una cartera diversificada). Se arma una cartera $W = 0.5X + 0.5Y$. Calcular $\text{Var}(W)$ y comparar con el promedio simple de las varianzas individuales.
\end{ejercicio}

\begin{solucion}
Usamos la fórmula general de varianza de una combinación lineal:
\[
\text{Var}(W) = \text{Var}(0.5X+0.5Y) = 0.5^2\text{Var}(X) + 0.5^2\text{Var}(Y) + 2(0.5)(0.5)\text{Cov}(X,Y).
\]
\[
= 0.25(100) + 0.25(25) + 0.5(-30) = 25 + 6.25 - 15 = 16.25.
\]
El promedio simple de las varianzas individuales sería $(100+25)/2 = 62.5$. La varianza real de la cartera ($16.25$) es mucho menor gracias a la covarianza negativa: diversificar entre activos que se mueven en sentidos opuestos reduce el riesgo total, ilustrando el principio básico de la teoría de carteras.
\end{solucion}

\begin{ejercicio}
Calcular la función generatriz de momentos de $X\sim\text{Bernoulli}(p)$ y usarla para obtener $E(X)$ y $\text{Var}(X)$.
\end{ejercicio}

\begin{solucion}
\[
M_X(t) = E(e^{tX}) = e^{t\cdot 0}(1-p) + e^{t\cdot 1}p = (1-p) + pe^t.
\]
Derivando: $M_X'(t) = pe^t$, por lo tanto $E(X) = M_X'(0) = pe^0 = p$.

Segunda derivada: $M_X''(t) = pe^t$, entonces $E(X^2)=M_X''(0)=p$.
\[
\text{Var}(X) = E(X^2)-[E(X)]^2 = p - p^2 = p(1-p).
\]
Resultado consistente con lo esperado para una Bernoulli.
\end{solucion}

\begin{ejercicio}
Sea $X\sim\text{Exp}(\lambda)$ con función generatriz de momentos $M_X(t)=\dfrac{\lambda}{\lambda-t}$ para $t<\lambda$. Usarla para obtener $E(X)$ y $E(X^2)$.
\end{ejercicio}

\begin{solucion}
Escribimos $M_X(t) = \lambda(\lambda-t)^{-1}$. Derivando:
\[
M_X'(t) = \lambda(\lambda-t)^{-2}, \qquad M_X'(0) = \lambda \cdot \lambda^{-2} = \frac{1}{\lambda} = E(X).
\]
\[
M_X''(t) = 2\lambda(\lambda-t)^{-3}, \qquad M_X''(0) = 2\lambda\cdot\lambda^{-3} = \frac{2}{\lambda^2} = E(X^2).
\]
Por lo tanto $\text{Var}(X) = \dfrac{2}{\lambda^2}-\dfrac{1}{\lambda^2} = \dfrac{1}{\lambda^2}$, resultado consistente con lo visto en la teoría.
\end{solucion}

\begin{ejercicio}
Sean $X_1\sim\text{Poisson}(3)$ y $X_2\sim\text{Poisson}(5)$ independientes. Usando la función generatriz de momentos, determinar la distribución de $Y = X_1+X_2$.
\end{ejercicio}

\begin{solucion}
La FGM de una Poisson$(\lambda)$ es $M(t) = e^{\lambda(e^t-1)}$. Como $X_1, X_2$ son independientes, por la propiedad de la FGM de una suma:
\[
M_Y(t) = M_{X_1}(t)\cdot M_{X_2}(t) = e^{3(e^t-1)}\cdot e^{5(e^t-1)} = e^{(3+5)(e^t-1)} = e^{8(e^t-1)}.
\]
Esta es exactamente la FGM de una $\text{Poisson}(8)$. Por la propiedad de \textbf{unicidad} de la FGM, concluimos
\[
Y = X_1+X_2 \sim \text{Poisson}(8).
\]
\end{solucion}

\begin{ejercicio}
Sean $X_1,X_2,X_3$ independientes con $X_i\sim N(\mu_i,\sigma_i^2)$, $\mu_1=10,\sigma_1^2=4$; $\mu_2=15,\sigma_2^2=9$; $\mu_3=5,\sigma_3^2=1$. Sea $T=X_1+X_2+X_3$. Hallar la distribución de $T$ usando la FGM, y calcular $P(T>32)$.
\end{ejercicio}

\begin{solucion}
La FGM de $N(\mu,\sigma^2)$ es $M(t)=e^{\mu t+\sigma^2t^2/2}$. Por independencia:
\[
M_T(t) = \prod_{i=1}^3 e^{\mu_i t+\sigma_i^2t^2/2} = e^{(\mu_1+\mu_2+\mu_3)t + (\sigma_1^2+\sigma_2^2+\sigma_3^2)t^2/2} = e^{30t+14t^2/2}.
\]
Esta es la FGM de una $N(30,14)$. Por unicidad, $T\sim N(30,14)$, con $\sigma_T=\sqrt{14}\approx 3.742$.

\medskip
Estandarizando:
\[
P(T>32) = P\left(Z > \frac{32-30}{3.742}\right) = P(Z>0.5345) \approx 1-\Phi(0.53) \approx 1-0.7019 = 0.2981.
\]
La probabilidad de que $T$ supere $32$ es aproximadamente $29.8\%$.
\end{solucion}

\begin{ejercicio}
Demostrar, usando la definición de covarianza, que si $X$ e $Y$ son independientes entonces $\text{Cov}(X,Y)=0$. (No es necesario suponer discretitud ni continuidad; puede usarse la propiedad $E(XY)=E(X)E(Y)$ para variables independientes.)
\end{ejercicio}

\begin{solucion}
Si $X$ e $Y$ son independientes, sabemos (propiedad general vista al definir independencia de variables aleatorias) que $E(XY) = E(X)\,E(Y)$.

\medskip
Por la fórmula de cálculo de la covarianza:
\[
\text{Cov}(X,Y) = E(XY) - E(X)E(Y) = E(X)E(Y) - E(X)E(Y) = 0. \qquad \blacksquare
\]
\medskip
\textbf{Observación:} la recíproca es falsa en general (ver el ejemplo de $Y=X^2$ con $X\sim U(-1,1)$ visto en la Clase 16): covarianza cero no implica independencia, solo ausencia de asociación \emph{lineal}.
\end{solucion}

\end{document}
